\section{Biblical passages for study}
The English in this section is the public-domain Douay--Rheims Bible, Challoner revision, in the Project Gutenberg 1581 witness. It is \textbf{not the text proclaimed in the celebration}; the U.S. \work{Lectionary for Mass}, no. 142, controls that English. Psalm numbers in these study texts follow the Vulgate. Whole-verse biblical bases below are not reconstructions of the adapted Missal antiphons or the Gospel acclamation. The approved English of the three orations is not reproduced; their identifiers and subjects appear in the inventory and commentary.

\subsection{Entrance: biblical basis}
\properlocus{Psalm 129:3--4 (Hebrew numbering 130). Complete study verses; the liturgical antiphon has its own form.}


\textsuperscript{3} If thou, O Lord, wilt mark iniquities: Lord, who shall stand it.

\textsuperscript{4} For with thee there is merciful forgiveness: and by reason of thy law, I have waited for thee, O Lord. My soul hath relied on his word:

\subsection{First reading}\properlocus{Isaiah 25:6--10a. The first clause of verse 10 alone concludes the appointment.}

\textsuperscript{6} And the Lord of hosts shall make unto all people in this mountain, a feast of fat things, a feast of wine, of fat things full of marrow, of wine purified from the lees.

\textsuperscript{7} And he shall destroy in this mountain the face of the bond with which all people were tied, and the web that he began over all nations.

\textsuperscript{8} He shall cast death down headlong for ever: and the Lord God shall wipe away tears from every face, and the reproach of his people he shall take away from off the whole earth: for the Lord hath spoken it.

\textsuperscript{9} And they shall say in that day: Lo, this is our God, we have waited for him, and he will save us: this is the Lord, we have patiently waited for him, we shall rejoice and be joyful in his salvation.

\textsuperscript{10} For the hand of the Lord shall rest in this mountain:

\subsection{Responsorial Psalm}\properlocus{Psalm 22:1--6 (Hebrew numbering 23); complete study psalm. The appointed response is 6cd in the Lectionary, not this complete final verse.}

\textsuperscript{1} A psalm for David. The Lord ruleth me: and I shall want nothing.

\textsuperscript{2} He hath set me in a place of pasture. He hath brought me up, on the water of refreshment:

\textsuperscript{3} He hath converted my soul. He hath led me on the paths of justice, for his own name’s sake.

\textsuperscript{4} For though I should walk in the midst of the shadow of death, I will fear no evils, for thou art with me. Thy rod and thy staff, they have comforted me.

\textsuperscript{5} Thou hast prepared a table before me against them that afflict me. Thou hast anointed my head with oil; and my chalice which inebriateth me, how goodly is it!

\textsuperscript{6} And thy mercy will follow me all the days of my life. And that I may dwell in the house of the Lord unto length of days.

\subsection{Second reading}\properlocus{Philippians 4:12--14, 19--20.}

\textsuperscript{12} I know both how to be brought low, and I know how to abound (every where and in all things I am instructed): both to be full and to be hungry: both to abound and to suffer need.

\textsuperscript{13} I can do all things in him who strengtheneth me.

\textsuperscript{14} Nevertheless, you have done well in communicating to my tribulation.

\notread{Verses 15--18 are not appointed.}

\textsuperscript{19} And may my God supply all your want, according to his riches in glory in Christ Jesus.

\textsuperscript{20} Now to God and our Father be glory, world without end. Amen.

\subsection{Acclamation: biblical basis}\properlocus{Ephesians 1:17--18. Complete study verses; the appointed acclamation is an adaptation (cf.), not these verses in full.}

\textsuperscript{17} That the God of our Lord Jesus Christ, the Father of glory, may give unto you the spirit of wisdom and of revelation, in the knowledge of him:

\textsuperscript{18} The eyes of your heart enlightened that you may know what the hope is of his calling and what are the riches of the glory of his inheritance in the saints.

\subsection{Gospel: both forms}\properlocus{Matthew 22:1--14; the short form ends after verse 10. The Douay opening remains that of its biblical witness.}

\textsuperscript{1} And Jesus answering, spoke again in parables to them, saying:

\textsuperscript{2} The kingdom of heaven is likened to a king who made a marriage for his son.

\textsuperscript{3} And he sent his servants to call them that were invited to the marriage: and they would not come.

\textsuperscript{4} Again he sent other servants, saying: Tell them that were invited, Behold, I have prepared my dinner: my beeves and fatlings are killed, and all things are ready. Come ye to the marriage.

\textsuperscript{5} But they neglected and went their ways, one to his farm and another to his merchandise.

\textsuperscript{6} And the rest laid hands on his servants and, having treated them contumeliously, put them to death.

\textsuperscript{7} But when the king had heard of it, he was angry: and sending his armies, he destroyed those murderers and burnt their city.

\textsuperscript{8} Then he saith to his servants: The marriage indeed is ready; but they that were invited were not worthy.

\textsuperscript{9} Go ye therefore into the highways; and as many as you shall find, call to the marriage.

\textsuperscript{10} And his servants going forth into the ways, gathered together all that they found, both bad and good: and the marriage was filled with guests.

\textbf{The short form ends here. Only the long form continues:}

\textsuperscript{11} And the king went in to see the guests: and he saw there a man who had not on a wedding garment.

\textsuperscript{12} And he saith to him: Friend, how camest thou in hither not having on a wedding garment? But he was silent.

\textsuperscript{13} Then the king said to the waiters: Bind his hands and feet, and cast him into the exterior darkness. There shall be weeping and gnashing of teeth.

\textsuperscript{14} For many are called, but few are chosen.

\subsection{Communion: psalm option, biblical basis}\properlocus{Psalm 33:11 (Hebrew numbering 34:11). The Missal option is an adaptation (cf.).}

\textsuperscript{11} The rich have wanted, and have suffered hunger: but they that seek the Lord shall not be deprived of any good.

\subsection{Communion: Johannine option, biblical basis}\properlocus{1 John 3:2. Complete study verse; the Missal excerpts its promise. This is an alternative to the psalm option.}

\textsuperscript{2} Dearly beloved, we are now the sons of God: and it hath not yet appeared what we shall be. We know that when he shall appear we shall be like to him: because we shall see him as he is.
